\clearpage
\section{Specyfika danych tabelarycznych i stan badań}
\label{chap:stan-badan}

\subsection{Definicja zadania i właściwości danych}
\label{sec:definicja-tabel}

Rozpatrywane zadanie opisuje zbiór $\mathcal{D}=\{(x_i,y_i)\}_{i=1}^{n}$, gdzie $x_i\in\mathcal{X}_1\times\cdots\times\mathcal{X}_d$ jest obserwacją o~ustalonym schemacie cech, a~$y_i\in\{1,\ldots,K\}$. Klasyfikator zwraca wektor prawdopodobieństw $p_\theta(y\mid x)$. Oznaczenia $n$, $d$ i~$K$ odnoszą się odpowiednio do liczby obserwacji, cech i~klas; wymiar reprezentacji po kodowaniu może różnić się od~$d$.

Specyfika tabel nie sprowadza się do współwystępowania zmiennych numerycznych i~kategorycznych. Poszczególne kolumny mają odrębne znaczenie, rozkłady i~mechanizmy powstawania braków, a~wybór kodowania określa geometrię wejścia dostępną modelowi~\cite{borisov2022survey}. Traktowanie kodu kategorii jako wielkości porządkowej wprowadza dodatkowe założenie o~relacji między wartościami. W~badaniu znaczenie tego założenia ujawnia się szczególnie w~zestawieniu metod odległościowych, sieci z~osadzeniami i~drzew z~natywną obsługą kategorii.

Odrębną decyzją jest jednostka predykcji. W~Diabetes wiersz opisuje pobyt, podczas gdy rozłączność treningu i~testu ma dotyczyć pacjentów. Identyfikator osoby określa zatem konstrukcję oceny, a~nie reprezentację predykcyjną. Takie rozróżnienie między cechą a~metadaną jest częścią definicji problemu. Tabela~\ref{tab:role-kolumn} zestawia jego konsekwencje dla przygotowania danych.

\begin{table}[!htbp]
\centering
\caption{Typy i role kolumn istotne w projektowaniu klasyfikacji. Zestawienie pojęciowe, nie opis konkretnego zbioru.}
\label{tab:role-kolumn}
\small
\begin{tabularx}{\textwidth}{@{}P{0.22\textwidth}YY@{}}
\toprule
Typ lub rola & Decyzja reprezentacyjna & Konsekwencja dla oceny \\
\midrule
Cecha numeryczna & Skala, transformacja rozkładu, osadzenie & Porównanie obejmuje geometrię wejścia i model. \\
Cecha nominalna & Kodowanie, rzadkie i nieznane poziomy & Słownik i statystyki kategorii zależą od treningu. \\
Cecha porządkowa & Wykorzystanie kolejności lub osobnych poziomów & Odstępy między kodami są dodatkowym założeniem. \\
Identyfikator grupy & Metadana podziału, nie predyktor & Jednostka rozłączności może obejmować wiele wierszy. \\
Czas obserwacji & Moment dostępności informacji & Losowy podział nie identyfikuje jakości przyszłej predykcji. \\
Etykieta docelowa & Semantyka zdarzenia i mapowanie klas & Kod dodatni nie musi oznaczać zdarzenia rzadkiego. \\
\bottomrule
\end{tabularx}
\end{table}

\subsection{Reprezentacja danych i założenia indukcyjne}
\label{sec:reprezentacja-zalozenia}

\subsubsection{Znaczenie reprezentacji kolumn}

W~danych tabelarycznych naturalny układ współrzędnych niesie informację: poszczególne osie odpowiadają określonym pomiarom lub kategoriom. Przestawienie kolumn razem z~ich definicjami jest zmianą zapisu, natomiast obrót mieszający cechy numeryczne zmienia sposób udostępnienia ich znaczenia modelowi. Nie ma podstaw, aby te dwie operacje traktować jako równoważne. Podobnie sąsiedztwo kolumn nie uzasadnia lokalności ani translacyjnego współdzielenia parametrów znanego z~przetwarzania obrazów~\cite{borisov2022survey}.

Kodowanie zero-jedynkowe, osadzenia kategorii i~transformacje cech numerycznych nie są neutralnymi operacjami wejściowymi. Pierwsze zachowuje rozróżnialność kategorii kosztem szerokości macierzy, drugie pozwala uczyć ich relacje, lecz wymaga oszacowania dodatkowych parametrów. Transformacja skali zmienia odległości wykorzystywane przez KNN i~warunki optymalizacji sieci. Z~tego powodu w~porównaniu głównym ocenie podlegają całe potoki, z~przetwarzaniem właściwym danej metodzie. Wymuszenie identycznego kodowania oznaczałoby ocenę innego pytania: jakości modeli przy narzuconej reprezentacji.

\subsubsection{Informacyjność cech i nieregularne zależności}

Grinsztajn i~współautorzy~\cite{grinsztajn2022trees} powiązali przewagę drzew w~badanym reżimie z~trzema własnościami: odpornością na cechy nieinformatywne, zdolnością odtwarzania nieregularnych zależności oraz wykorzystaniem orientacji danych. W~ich interwencjach usuwanie mało informatywnych cech zmniejszało różnicę między sieciami typu MLP a~pozostałymi modelami, natomiast dodawanie cech losowych ją zwiększało. Obrót numerycznego wejścia zmieniał względną kolejność metod, mimo zachowania informacji zawartej w~obserwacjach. Wyniki dotyczą konkretnych architektur i~przekształceń, ale pokazują, dlaczego sama liczba parametrów nie wyjaśnia zdolności uczenia na tabelach.

Z~perspektywy tej pracy prowadzi to do rozdzielenia dwóch problemów: elastyczności reprezentacji i~efektywnego wykorzystania dostępnej próby. Szeroki zbiór nie musi być trudny, jeżeli sygnał jest silny, a~użyteczna zależność prosta. Z~kolei zwiększanie pojemności sieci nie gwarantuje poprawy, gdy dodatkowe stopnie swobody służą dopasowaniu cech nieinformatywnych. Bioresponse stanowi przypadek wysokiej wymiarowości względem liczności treningu, lecz bez kontrolowanych rotacji i~dodawania cech nie pozwala przypisać różnic jakości samemu układowi współrzędnych. Takie rozstrzygnięcie wymagałoby osobnej interwencji.

\subsubsection{Braki i informacje niedostępne przy predykcji}

Brak wartości nie jest tym samym co zero ani nowa kategoria występująca po raz pierwszy podczas predykcji. Procedura musi określać, jak te sytuacje reprezentuje. Jeżeli imputacja lub kodowanie uczą się statystyk, ich dopasowanie powinno respektować granicę między treningiem i~oceną. Włączenie etykiety do kodowania kategorii wymaga szczególnej ostrożności, ponieważ statystyka może ujawnić informację, której model nie otrzyma dla nowego przypadku~\cite{prokhorenkova2018catboost}.

Poprawna obsługa techniczna braków nie wystarcza do stwierdzenia poprawności zadania. Cecha zarejestrowana dopiero po wystąpieniu przewidywanego zdarzenia może zapewniać wysoką jakość bez użyteczności w~zadanym momencie predykcji. Dostępność informacji i~role cech określają dopuszczalną interpretację predykcji; nie wynikają z~samej jakości modelu.

\subsection{Trudności w zadaniach klasyfikacji}
\label{sec:trudne-rezimy}

\subsubsection{Mała liczba przykładów i wysoka wymiarowość}

Dla oceny generalizacji istotna jest liczba przykładów dostępnych w~konkretnym dopasowaniu, a~nie całkowity rozmiar pliku. Po wydzieleniu testu, walidacji i~podzbioru treningowego relacja wymiaru reprezentacji do liczby obserwacji może być znacznie mniej korzystna niż sugeruje opis zbioru. Dodatkowo rzadkie kategorie i~klasy mają własne, mniejsze liczności efektywne. Przy tej samej wartości $n$ dwa zadania mogą zatem dostarczać bardzo różnej ilości informacji o~parametrach decydujących o~predykcji.

Badanie Kadry i~współautorów~\cite{kadra2021regularization} na 40 zbiorach uzasadnia uwzględnienie wspólnie dobieranych technik regularyzacji jako istotnego składnika porównania sieci. Taki wynik nie rozstrzyga jednak, czy dany składnik pomaga przy ustalonych pozostałych ustawieniach. W~badaniu własnym przyjęto właśnie tę drugą perspektywę: usunięcie dropoutu albo zaniku wag z~jednej procedury MLP, przy zmiennej liczności treningu. Przedmiotem oceny jest efekt warunkowy względem architektury, optymalizatora i~reguły zatrzymania, a~nie najlepszy osiągalny wynik każdej osobno strojonej rodziny.

Zmiana udziału treningu przy stałej liczbie epok może zmieniać również liczbę aktualizacji parametrów. Jest to istotne dla interpretacji zarówno optymalizacji, jak i~skumulowanego zaniku wag. Dlatego krzywe jakości względem liczności są interpretowane łącznie z~przebiegiem uczenia i~wybraną epoką. Stałe walidacja i~test oddzielają zmianę treningu od zmiany ocenianych populacji. Taka zależność stanowi krzywą wrażliwości na liczność treningu, przy zachowanym budżecie etykiet walidacyjnych.

\subsubsection{Niezrównoważenie i rozpoznawanie rzadkiej klasy}

W~klasyfikacji niezrównoważonej trzeba rozdzielić trzy obiekty oceny: uporządkowanie obserwacji, wartości prawdopodobieństw i~decyzje przy zadanym progu. ROC AUC opisuje pierwszy z~nich, ale nie określa udziału prawdziwych zdarzeń wśród alarmów. Przy częstości klasy rzadkiej $\pi$ precyzja dla ustalonego progu spełnia tożsamość
\begin{equation}
\operatorname{Precision}=\frac{\pi\operatorname{TPR}}
{\pi\operatorname{TPR}+(1-\pi)\operatorname{FPR}},
\label{eq:precyzja-czestosc}
\end{equation}
gdzie TPR oznacza czułość, a~FPR odsetek fałszywych alarmów w~klasie pozostałej. Zależność wynika bezpośrednio z~definicji prawdopodobieństwa warunkowego: przy dodatnich, niezmienionych TPR i~FPR spadek $\pi$ obniża precyzję. Uzasadnia to analizę precision--recall obok ROC, zgodnie z~argumentacją Saito i~Rehmsmeiera~\cite{saito2015pr}. W~badaniu miarą oceny jest average precision (AP), liczona jako suma przyrostów czułości ważonych precyzją, a~nie jako trapezowa całka krzywej.

Rozróżnienie ważenia i~progowania można doprecyzować analitycznie. Niech $\eta(x)=P(Y=1\mid X=x)$ oznacza prawdopodobieństwo rzadkiego zdarzenia, przy czym kod~1 przyjęto wyłącznie na potrzeby poniższego wyprowadzenia. Dla dodatnich, stałych wag klas $w_1,w_0$ warunkowe ryzyko ważonej straty logarytmicznej ma postać
\begin{equation}
R_w(q\mid x)=-w_1\eta(x)\log q-w_0[1-\eta(x)]\log(1-q).
\label{eq:ryzyko-wazone}
\end{equation}
Wyzerowanie pochodnej po $q$ daje, dla $0<\eta(x)<1$,
\begin{equation}
q_w^*(x)=\frac{w_1\eta(x)}{w_1\eta(x)+w_0[1-\eta(x)]},
\qquad
\operatorname{logit}q_w^*(x)=\operatorname{logit}\eta(x)+\log\frac{w_1}{w_0}.
\label{eq:optimum-wazenia}
\end{equation}
Powyższe wyprowadzenie dotyczy nieograniczonej klasy funkcji i~ryzyka populacyjnego, nie rzeczywistego przebiegu optymalizacji. W~tych warunkach ważenie zmienia skalę prawdopodobieństw, lecz zachowuje uporządkowanie: transformacja jest ściśle rosnąca. Próg $q_w^*\geq 1/2$ odpowiada progowi $\eta\geq w_0/(w_0+w_1)$. Sam wzrost czułości po ważeniu nie dowodzi więc lepszej reprezentacji ani lepszego rankingu.

W~skończonej próbie, przy ograniczonej pojemności i~regularyzacji, ważenie może zmienić również wyuczone uporządkowanie. Dlatego badanie~A obejmuje równolegle AP, ROC AUC, nieważoną stratę logarytmiczną oraz decyzje przy progach wybranych na walidacji. Pozwala to odróżnić poprawę rankingu od przesunięcia punktu pracy i~zmiany jakości prawdopodobieństw. Zdarzenie rzadkie jest identyfikowane semantycznie, niezależnie od kodu nadanego mu przez procedurę przygotowania danych.

\subsubsection{Wieloklasowość i jakość prawdopodobieństw}

W~zadaniu wieloklasowym strata logarytmiczna ocenia cały rozkład predykcyjny poprzez prawdopodobieństwo przypisane etykiecie rzeczywistej. Dla ustalonego $x$, prawdziwego rozkładu $\eta_c=P(Y=c\mid x)$ i~predykcji $p_c$ jej wartość oczekiwaną można zapisać jako
\begin{equation}
-\sum_{c=1}^{K}\eta_c\log p_c
=H(\eta)+D_{\mathrm{KL}}(\eta\Vert p).
\label{eq:logloss-kl}
\end{equation}
Dekompozycja wynika z~dodania i~odjęcia $\sum_c\eta_c\log\eta_c$. Przy dostępie do prawdziwego rozkładu minimum osiąga $p=\eta$; wysoka pewność błędnej klasy jest karana nawet wtedy, gdy liczba poprawnych decyzji pozostaje bez zmian. Niższa strata na skończonym teście nie identyfikuje jednak osobno poprawy kalibracji: zależy również od zdolności rozróżniania przypadków.

W~przyjętym protokole strata logarytmiczna jest główną miarą zadań wieloklasowych, zgodnie z~adaptacją TabArena~\cite{erickson2025tabarena}. Covertype, Diabetes i~Helena reprezentują różne liczby klas. W~Helenie istotna jest nie tylko obsługa stu wyjść, lecz także rozkład liczności klas i~jakość oszacowania prawdopodobieństw przy ograniczonej liczbie przykładów niektórych etykiet. Wartości straty są porównywane między metodami w~obrębie zadania; ich bezpośrednie uśrednienie między zadaniami mieszałoby różne poziomy trudności i~entropii etykiet.

\subsubsection{Zależności między obserwacjami i zmiana rozkładu}

Powtarzające się obserwacje tej samej osoby mogą naruszać interpretację losowego podziału jako oceny nowych osób. Podobnie losowe przemieszanie danych z~różnych okresów nie musi odpowiadać przewidywaniu przyszłości. Rubachev i~współautorzy~\cite{rubachev2025tabred} wskazali znaczenie podziałów czasowych w~TabReD, obejmującym osiem przemysłowych zbiorów z~rozbudowanymi cechami. Zmiana sposobu podziału wpływała na względną ocenę metod.

W~tej pracy podział Diabetes zachowuje rozłączność pacjentów. Zakres badania nie obejmuje osobnego porównania czasowego, dlatego nie uzasadnia wniosków o~odporności na dryf. Dla pozostałych źródeł interpretacja wyniku również zależy od tego, czy podział odzwierciedla właściwą populację predykcji.

\subsection{Rozwój podejść neuronowych}
\label{sec:rozwoj-podejsc}

Rozwiązania neuronowe można porządkować według tego, na którym poziomie ingerują w~procedurę predykcji. Zmiana reprezentacji przygotowuje wejście dla modelu, zmiana architektury określa sposób łączenia informacji, a~procedura uczenia wpływa na odnajdywane parametry. Podział ten jest użyteczny, lecz nie jest rozłączny: jedna metoda może łączyć kilka takich elementów~\cite{borisov2022survey}.

Przyjęty dobór obejmuje także podejście oparte na sąsiedztwie i~model wykorzystujący wiedzę nabytą przed otrzymaniem badanego zadania. Ich obecność rozszerza porównanie poza klasyczny schemat trenowania niezależnej sieci od początku na każdym zbiorze. Podział na modele neuronowe i~konwencjonalne pozostaje przydatny organizacyjnie, ale nie wyczerpuje różnic między procedurami.

Rysunek~\ref{fig:mapa-podejsc} porządkuje wybrane metody według ich roli w~porównaniu. Wyjaśnienie mechanizmów znajduje się w~rozdziale~\ref{chap:metody}.

\plannedfigure
    {figures/ch02/mapa-podejsc.pdf}
    {Role wybranych podejść w~badaniu. Opracowanie własne.}
    {fig:mapa-podejsc}
    {\begin{tabular}{ll}
Odniesienie gęste & MLP \\
Selekcja cech i~uwaga & TabNet, FT-Transformer \\
Procedura uczenia & RealMLP \\
Uczone sąsiedztwo & ModernNCA \\
Współdzielenie parametrów zespołu & TabM \\
Pretrenowanie i~kontekst & TabPFN-3.5 \\
Odniesienia konwencjonalne & Regresja logistyczna, KNN, XGBoost, CatBoost
\end{tabular}}

\subsection{Badania porównawcze}
\label{sec:benchmarki}

Przegląd ma charakter narracyjny i~obejmuje publikacje uwzględnione przy doborze metod do 27 października 2026 roku. Dla preprintów wskazano wykorzystane wersje, ponieważ kolejne rewizje mogą różnić się zakresem badań. Przegląd nie był systematycznym przeszukaniem całej literatury dziedziny.

\subsubsection{Od porównania architektur do analizy warunków}

Gorishniy i~współautorzy~\cite{gorishniy2021revisiting} porównali architektury na jedenastu zadaniach klasyfikacji i~regresji we wspólnej procedurze uczenia i~strojenia. Istotnym wynikiem było wykazanie konkurencyjności prostego ResNet oraz silnej pozycji FT-Transformer wśród badanych sieci. Porównanie z~drzewami nie wskazało rozwiązania wygrywającego we wszystkich zadaniach. Dla projektu tej pracy oznacza to konieczność uwzględnienia mocnego odniesienia gęstego: specjalizacja architektury wymaga wykazania korzyści względem właściwie przygotowanej sieci bazowej.

Grinsztajn i~współautorzy~\cite{grinsztajn2022trees} uzyskali przewagę drzew na benchmarku 45 zbiorów, utrzymującą się wraz ze zwiększaniem liczby prób strojenia. Badanie obejmowało jednak określony reżim: ograniczoną liczność treningu, usunięte braki, równoważenie klasyfikacji oraz ograniczenia wymiarowości i~kategorii. Wynik stanowi argument za zachowaniem silnych odniesień drzewiastych, lecz nie rozstrzyga o~naturalnie rzadkich klasach ani o~wszystkich późniejszych procedurach neuronowych.

McElfresh i~współautorzy~\cite{mcelfresh2024tabzilla} rozszerzyli perspektywę do 19 algorytmów i~176 zbiorów klasyfikacji. W~ich analizie na około jednej trzeciej zadań przyrost jakości po lekkim strojeniu CatBoost był większy niż różnica między najlepszym domyślnym modelem drzewiastym i~neuronowym. Znaczenie ma tu konkretna konstrukcja porównania, a~nie ogólne pierwszeństwo strojenia przed doborem algorytmu. Dla niniejszej pracy wynik uzasadnia oddzielną ocenę ustawień domyślnych i~strojonych. Wczesny TabPFN był przy tym oceniany w~odmiennych warunkach i~nie stanowi dowodu skuteczności wybranej tutaj wersji~3.5.

\subsubsection{Procedury uczenia i rozwój benchmarków}

Holzmüller i~współautorzy~\cite{holzmueller2024realmlp} opracowali ustawienia RealMLP na 118 zbiorach, a~następnie ocenili ich przenoszenie na 90 zbiorów meta-test. Uzyskali konkurencyjność rozwiniętej procedury gęstej względem drzew, lecz obraz zależał od kolekcji i~metryki: na benchmarku Grinsztajna RealMLP wypadał słabiej od CatBoost w~klasyfikacji. Ważnym ustaleniem jest zatem skuteczność przenoszenia ustawień między zadaniami, a~nie bezwarunkowe rozstrzygnięcie na korzyść MLP. Konfiguracja domyślna może zawierać rezultat kosztownego wcześniejszego rozwoju, nawet jeżeli nie jest strojona na aktualnym zbiorze.

Wersja~4 pracy Ye i~współautorów~\cite{ye2025talent} obejmuje 40 metod na 300 bazowych zbiorach oraz dodatkowe kolekcje. W~zestawieniach wysokie pozycje zajmują modele pretrenowane, a~wśród metod uczonych zadaniowo wyróżniają się rozwinięte MLP i~ModernNCA. Starsze specjalizacje, w~tym TabNet, wypadają słabiej od czołowych zespołów. Autorzy odnotowują także korzyści zespołowania, między innymi dla TabM. To uzasadnia reprezentowanie kilku kierunków rozwoju sieci zamiast utożsamiania całej rodziny z~jednym MLP. Główny protokół stosuje podział 64/16/20 i~100 prób strojenia; agregaty obejmujące regresję pozostają odrębne od wniosków o~klasyfikacji.

TabArena~\cite{erickson2025tabarena} przenosi punkt ciężkości z~jednorazowego rankingu na wersjonowany system oceny danych, modeli i~procedur wyszukiwania. Początkowa kolekcja obejmuje 51 zbiorów. Dla niniejszego projektu najistotniejsze jest rozdzielenie konfiguracji domyślnej, wyboru po strojeniu i~zespołu konfiguracji. Pozwala ono ocenić, czy poprawa wynika z~dodatkowego wyszukiwania, czy z~wykorzystania różnorodności znalezionych rozwiązań. Przyjęte dwanaście zadań i~własne wyjątki walidacyjne tworzą adaptację tej organizacji eksperymentu, nie reprodukcję oficjalnego rankingu.

\begin{table}[tbp]
\centering
\caption{Wybrane badania porównawcze i ich znaczenie dla projektu pracy. Liczby dotyczą wskazanych publikacji lub wersji, a nie jednego wspólnego benchmarku. Opracowanie własne na podstawie cytowanych źródeł.}
\label{tab:benchmarki}
\footnotesize
\begin{tabularx}{\textwidth}{@{}P{0.25\textwidth}P{0.22\textwidth}Y@{}}
\toprule
Praca & Zakres & Ograniczenia porównania \\
\midrule
RTDL~\cite{gorishniy2021revisiting} & 11 zadań & Klasyfikacja i regresja; starsze modele i inny protokół. \\
Grinsztajn i in.~\cite{grinsztajn2022trees} & 45 zbiorów & Kontrolowany reżim nie zachowuje wszystkich trudności danych źródłowych. \\
TabZilla~\cite{mcelfresh2024tabzilla} & 176 zbiorów klasyfikacji & Poszczególne analizy obejmują różne podzbiory zadań i ustawienia modeli. \\
RealMLP~\cite{holzmueller2024realmlp} & 118 meta-train, 90 meta-test & Rozwój procedury jest odrębny od oceny i kosztów zadaniowych. \\
TALENT v4~\cite{ye2025talent} & 300 zbiorów bazowych i rozszerzenia & Mieszane typy zadań i inna walidacja. \\
TabArena~\cite{erickson2025tabarena} & 51 zbiorów początkowych & Zespoły, wersje i zasoby są częścią ocenianej procedury. \\
TabReD~\cite{rubachev2025tabred} & 8 zbiorów przemysłowych & Podziały czasowe nie odpowiadają standardowemu IID. \\
\bottomrule
\end{tabularx}
\end{table}

\subsection{Zgodności, rozbieżności i ograniczenia syntezy}
\label{sec:synteza-literatury}

Z~zestawienia wynika przede wszystkim, że etykieta rodziny modelu jest zbyt zgrubna, aby wyjaśniać postęp. Prace RTDL i~RealMLP podważają ocenę nowych architektur względem słabo skonfigurowanej sieci gęstej; TALENT pokazuje znaczenie rozwiniętych procedur i~zespołów. Wspólnym wnioskiem przyjętym w~tej pracy jest konieczność porównania konkretnych konfiguracji wraz z~ich reprezentacją i~kosztem selekcji, a~nie abstrakcyjnych kategorii ,,sieci'' i~,,drzewa''~\cite{gorishniy2021revisiting,holzmueller2024realmlp,ye2025talent}.

Rozbieżność między przewagą drzew w~badaniu Grinsztajna a~konkurencyjnością nowszych sieci nie jest prostym odwróceniem jednego wyniku. Zmieniono zarówno procedury, jak i~populacje zadań, budżety oraz kryteria agregacji. RealMLP dostarcza szczególnie użytecznego przykładu: jego względna pozycja zmienia się między kolekcjami ocenianymi w~tej samej publikacji~\cite{grinsztajn2022trees,holzmueller2024realmlp}. W~niniejszym badaniu uzasadnia to prezentację różnic zadaniowych obok rang. Nie przeprowadzono metaanalizy wyznaczającej wspólną wielkość przewagi z~tych niejednorodnych protokołów.

Pozostaje również różnica między wyjaśnieniem wyniku a~wyborem skutecznej procedury. Benchmark ocenia efekt całego rozwiązania; interwencja usuwa lub zmienia wskazany składnik. Wyniki tych dwóch rodzajów badań odpowiadają na uzupełniające się pytania. Powiązanie rankingu z~metacechami może sugerować hipotezę, lecz dopiero dopasowana kontrola pozwala odnieść różnicę do zadanej zmiany. Właśnie na tej granicy oparto podział własnego badania na porównanie główne i~dwa eksperymenty uzupełniające.

Ważną osią porównywalności jest oddzielenie selekcji od oceny. Cawley i~Talbot~\cite{cawley2010selection} pokazują, że dopasowywanie się do skończonej próby walidacyjnej może obciążać ocenę wybranej procedury. Wspólna walidacja nie daje zatem prawa do wielokrotnego projektowania zmian na podstawie testu. Etykiety testowe nie powinny służyć wyborowi konfiguracji, progu, zespołu ani zakresu prezentowanych porównań.

Drugą osią jest jednostka wnioskowania. Porównanie między zbiorami różni się od analizy powtórzeń na jednym zbiorze~\cite{demsar2006comparisons}. W~niniejszym badaniu foldy i~konfiguracje nie są traktowane jako nowe niezależne zadania. Rangi uzupełniają wyniki zadaniowe, ponieważ pomijają wielkość różnicy. Ponadto brak informacji o~pokryciu danych rozwoju metod nie jest zastępowany założeniem, że wszystkie wybrane zadania są całkowicie nowe dla ich twórców.

\subsection{Uzasadnienie własnego projektu}
\label{sec:uzasadnienie-projektu}

Własne badanie łączy ocenę użyteczności kompletnych procedur z~analizą dwóch mechanizmów istotnych dla trudnych zadań klasyfikacji. Porównanie główne odpowie, jaką jakość i~koszt uzyskuje się przy określonym sposobie selekcji modelu. Nie wyjaśni jednak, czy ewentualny wzrost wykrywalności rzadkiej klasy wynika z~lepszego rankingu, ani czy zachowanie MLP przy małej próbie zależy od danego składnika regularyzacji. Temu służą ważenie klas oraz dopasowane ablacje przy zagnieżdżonych podzbiorach treningowych.

Wartość tego układu polega na wspólnej interpretacji jakości, kosztu i~reakcji na kontrolowaną zmianę. Nie zakłada się, że każdy wybrany model jest skuteczny w~każdym reżimie. Możliwym wynikiem jest zarówno korzyść mechanizmu, jak i~jej brak lub odwrócenie po zmniejszeniu treningu. Tabela~\ref{tab:problem-dowod} określa, które pytania mają bezpośrednie pokrycie w~projekcie eksperymentu, a~które stanowią kontekst interpretacji.

\begin{table}[tbp]
\centering
\caption{Powiązanie właściwości danych z~zakresem badania.}
\label{tab:problem-dowod}
\small
\begin{tabularx}{\textwidth}{@{}P{0.26\textwidth}YY@{}}
\toprule
Problem & Zakres porównania & Granica interpretacji \\
\midrule
Jakość i koszt potoków & Wspólne testy, D/T/T+E & Bez izolacji pojedynczej architektury. \\
Rzadkie klasy & Ważenie i kontrola progów & Efekt przy ustalonych pozostałych ustawieniach. \\
Mało danych i regularyzacja & Zagnieżdżone podzbiory, ablacje MLP & Bez pełnej analizy interakcji mechanizmów. \\
Wymiarowość i typy cech & Zróżnicowane naturalne zadania & Związek opisowy, nie kontrolowany efekt jednej cechy. \\
Szum, rotacje, dryf & Omówienie literatury i audyt & Brak własnego testu odporności w tych reżimach. \\
\bottomrule
\end{tabularx}
\end{table}

Badanie główne pozwala porównać jakość i~koszt wybranych metod w~jednym protokole. Badania uzupełniające zawężają analizę do wpływu ważenia klas oraz dwóch składników regularyzacji MLP. Mechanizmy tych metod i~uzasadnienie ich wyboru przedstawiono w~kolejnym rozdziale.
